\vssub
\subsection{~嵌套} \label{sub:num_nest}
\conthead{\ws}{H. L. Tolman}

\noindent
传统上，波浪模式只考虑单向嵌套，即由低分辨率网格向高分辨率网格提供边界数据。\ws{} 一直提供这种方法，相关内容见 \para\ref{sub:one_way}。在模式 3.14 版中，引入了多网格波浪模式驱动程序，考虑网格之间的完整双向嵌套。该方法见 \para\ref{sub:two_way}。下文图示考虑规则网格，但所讨论的原则同样适用于曲线网格和三角形网格。


% -------------------------------------------------------------
\vssub
\subsubsection{~传统单向嵌套} \label{sub:one_way}
\vssub

\begin{figure}
\begin{center}

\setlength{\unitlength}{0.001in}

\begin{picture}(3000,1800)(0,-1550)

\put(0200,-1200){\vector(1,0){2100}}
\put(2400,-1250){\makebox(400,100){\small 空间}}
\put(0300,-1300){\vector(0,1){1300}}
\put(0100, 0100){\makebox(400,100)[c]{\small 时间}}

\multiput(0300,-1200)(0,300){4}{\circle {50}}
\multiput(0600,-1200)(0,300){4}{\circle {50}}
\multiput(0900,-1200)(0,300){4}{\circle {50}}
\multiput(1200,-1200)(0,300){4}{\circle*{50}}
\multiput(1500,-1200)(0,300){4}{\circle {50}}
\multiput(1800,-1200)(0,300){4}{\circle {50}}
\multiput(2100,-1200)(0,300){4}{\circle {50}}

\put(0300,-0150){\makebox(600,100)[c]{“陆地”}}
\put(0900,-0200){\vector(-1,0){575}}

\put(1500,-0150){\makebox(600,100)[c]{海洋}}
\put(1500,-0200){\vector(1,0){875}}
				       
\put(0900, 0050){\makebox(600,100)[c]{边界数据}}
\put(1100,-1200){\dashbox{50}(0200,1150){}}

\put(1050,-1550){\makebox(600,100)[c]{边界格式}}
\put(1200,-1400){\dashbox{25}(0300,1100){}}

\put(1550,-0750){\makebox(600,100)[l]{内部格式}}
\put(1500,-0800){\vector(1,0){875}}

\end{picture}
\end{center}

\caption{{\file ww3\_shel} 中使用的传统单向嵌套方法。空间和时间上的一维表示，符号表示网格点。} \label{fig:nest1}
\botline
\end{figure}


常规波浪模式程序 {\file ww3\_shel} 考虑单个波浪模式网格。该程序包含将边界条件从大尺度运行传递到小尺度运行的选项。每次运行可同时接受一个边界条件数据集，并生成最多 9 个边界条件数据集。为在水深和流场不兼容时保证波能守恒，边界数据由能量谱 $F(\sigma,\theta)$ 组成。数据文件包含生成运行中网格点处的谱，以及在所请求边界点处插值谱所需的信息。因此，如果小尺度运行的输入点位于大尺度运行的网格线上，传输文件的大小就最小。作为输入使用时，谱会在每个全局时间步长 $\Delta t_g$ 上以谱分量线性插值的方式在空间和时间上插值。

在波浪模式中包含边界数据的数值方法如图~\ref{fig:nest1} 所示。活动边界点在网格中被指定，用于将海点与陆点或其他停用网格点分隔开。在活动边界点和海点之间，应用局地边界格式（通常为一阶）。在模式内部海点中，则使用所选传播格式。

传统单向嵌套方法的实际使用方面在附录~\ref{app:nest} 中有更详细讨论。


% -------------------------------------------------------------
\vssub
\subsubsection{~双向嵌套} \label{sub:two_way}
\vssub

模式 3.14 版包含在程序 {\file ww3\_multi} 中使用多网格或镶嵌方法进行波浪建模的选项 \citep{tol:Vict06a,
  tol:MMAB07b, tol:OMOD08b}。在该程序中，可考虑任意数量、任意分辨率的网格，并在每个相关模式时间步长上在网格之间交换数据。网格被赋予等级编号，较低等级对应较低分辨率，相同等级对应相近分辨率（但不一定完全相同）。网格之间考虑三类数据传输：

\begin{itemize}
\item 从低等级网格向高等级网格传输数据。
\item 从高等级网格向低等级网格传输数据。
\item 在相同等级网格之间传输数据。
\end{itemize}

从低等级网格向高等级网格传输数据，是通过向高等级网格提供边界数据来完成的，方式与上一节和图~\ref{fig:nest1} 中所述的传统单向嵌套方法相同。

\begin{figure}
\begin{center}

\setlength{\unitlength}{0.001in}

\begin{picture}(3000,1900)(0,-1600)

\multiput(0300,-1400)(0,300){6}{\circle {50}}
\multiput(0600,-1400)(0,300){6}{\circle {50}}
\multiput(0900,-1400)(0,300){6}{\circle {50}}
\multiput(1200,-1400)(0,300){6}{\circle {50}}
\multiput(1500,-1400)(0,300){6}{\circle {50}}
\multiput(1800,-1400)(0,300){6}{\circle {50}}
\multiput(2100,-1400)(0,300){6}{\circle {50}}
\multiput(2400,-1400)(0,300){6}{\circle {50}}
\multiput(2700,-1400)(0,300){6}{\circle {50}}

\multiput(0450,-1550)(300,0){8}{\dashbox{50}(0000,1800){}}
\multiput(0150,-1250)(0,300){5}{\dashbox{50}(2700,0000){}}

\multiput(0400,-1450)(0,800){3}{\circle*{100}}
\multiput(1400,-1450)(0,800){3}{\circle*{100}}
\multiput(2400,-1450)(0,800){3}{\circle*{100}}

\put(895,-1055){\framebox(1010,810){}}
\put(900,-1050){\framebox(1000,800){}}
\put(905,-1045){\framebox(990,790){}}

\end{picture}
\end{center}

\caption{双向嵌套方法中协调低等级网格与高等级网格的概念示意。$\circ$ 和虚线分别表示高等级网格点和网格单元，$\bullet$ 和实线分别表示低等级网格及其中心网格单元。}
\label{fig:nest2}
\botline
\end{figure}


当该方法与从高等级到低等级的数据传输结合时，就建立了完整的双向嵌套方法。在 {\file ww3\_multi} 中，当高等级网格在时间上“追上”低等级网格之后，会将低等级网格中的数据与高等级网格中的数据相协调。考虑到低等级网格的分辨率按定义低于高等级网格的分辨率，从高等级网格中的能量 $E_{h,j}$ 估计低等级网格中的波能 $E_{l,i}$ 的一种自然方式为

\begin{equation}
E_{l,i} = \sum w_{i,j} E_{h,j} \:\:\: , \label{eq:nest_hg1}
\end{equation}

\noindent
其中 $i$ 和 $j$ 是两个网格中的网格计数器，$w_{i,j}$ 是平均权重。权重可按与波能守恒一致的方式定义为：高等级网格中覆盖低等级网格单元 $i$ 的网格单元 $j$ 的面积，并用低等级网格单元 $i$ 的面积归一化。图~\ref{fig:nest2} 对此作了说明。为避免循环协调，低等级网格中为高等级网格提供边界数据的网格点不以这种方式更新。

\input{zh/num/fig_nest3_zh}

相同等级的重叠网格不能使用上述双向嵌套技术来一致地交换数据。相反，所有这类网格都先传播一个时间步长，随后如图~\ref{fig:nest3} 所示对网格进行协调。对于网格 1（图~\ref{fig:nest3} 中的 $\circ$），可区分两个区域。在区域 C 中，自上次协调以来，边界影响已传播进入网格。实际穿透深度取决于数值格式的模板宽度以及传播时间步数。在区域 A 和 B 中，来自边界的信息尚未穿透，因此该区域可视为网格 1 的“内部”。类似地，区域 A 表示网格 2（图~\ref{fig:nest3} 中的 $\bullet$）的边界穿透深度，而 B 和 C 表示网格 2 的内部。网格 1 和网格 2 之间一种简单且一致的协调方法，是在区域 A 中仅使用网格 1 的数据（必要时将网格 1 的数据插值到网格 2 的网格点），并在区域 C 中仅使用网格 2 的数据。在区域 B 中，两个网格的内部部分重叠，可通过使用两个网格的加权平均获得一致解。注意，该方法很容易扩展到多个重叠网格。

注意，对于显式数值传播格式，以及具有相同分辨率且网格点重合的重叠网格，只要重叠区域足够宽，重叠网格的解就可以与兼容的单一网格解相同。

\vspace{\baselineskip} 
\noindent 
{\file ww3\_multi} 中的双向嵌套技术在很大程度上是自动化的。每个网格都独立准备，并具有自己首选的时间步进信息。每个网格期望从低等级网格获得边界数据的位置，会像单向嵌套方法一样标记。实现双向嵌套技术所需的所有其他簿记操作都是自动化的，不过可能需要若干次迭代，以确保每个网格中定义的所有输入边界点都能从多网格应用中的其他网格获得边界数据。作为替代，每个网格也可像传统嵌套方法一样从外部数据文件获取数据。在当前实现中，每个网格必须从单个文件或从多网格应用中的其他网格获得全部边界数据，但不能同时从文件和网格接收数据。关于为同时运行所有网格而开发的管理算法，详见 \citet[section 3.4]{tol:MMAB07b} 和 \cite{tol:OMOD08b}，此处不再重复。

注意，{\file ww3\_multi} 中使用的网格不需要具有相同的谱离散。谱会在 {\file
ww3\_multi} 中即时转换。该方法所用数值技术的细节可见 \citet[section 3.5.5]{tol:MMAB07b}。在这种方法中生成多网格可能很繁琐，并且需要保证网格之间的一致性以获得一致的模式结果。因此，\cite{tol:MMAB07a, tol:OMOD08a} 开发了自动网格生成工具。
